\section{Exact operator seeds: a mathematical contract}
\label{sec:exact-core}
This chapter defines the proposed \code{XOP-R1} operator-seed profile for
release 3.6.1.11. Definitions specify a representation and its meaning;
propositions establish conditional mathematical consequences. An implementation
must still discharge the stated parsing, arithmetic and proof-checking
obligations. No new runtime, machine-checked proof, test, simulation or measured
improvement is asserted. The implemented R10 orbital seed and its binary64
arithmetic keep their existing format and interpretation.

\subsection{Bits, values and exact expressions are distinct types}
Write $\mathsf{Bits}(w)=\B^w$ for a finite raw bit string. Its unsigned
interpretation is an explicitly requested map
\[
 \operatorname{uint}(b)=\sum_{i=0}^{w-1}b_i2^i.
\]
A binary64 payload belongs to $\mathsf{Bits}(64)$ until an explicit conversion
is applied. It does not become an arbitrary-precision real by being packed in
one-bit planes. A finite binary64 value can be converted to the exact dyadic
rational denoted by that value; this preserves its rounded input value, not an
unknown physical quantity. Signed zeros are distinct raw payloads but the same
rational zero. Infinity and NaN have no value in the exact real scalar types.

The scalar hierarchy is
\[
 \Z\ \lhook\joinrel\longrightarrow\ \mathbb Q
 \ \lhook\joinrel\longrightarrow\ \mathbb A
 \ \lhook\joinrel\longrightarrow\ \mathsf{RealTerm},
 \qquad \mathbb A=\overline{\mathbb Q}\cap\R.
\]
Here $\mathsf{RealTerm}$ is finite syntax for a possibly partial real expression,
not a claim that every real number has a finite description. Coercions preserve
denotation and are explicit in the typed graph. Booleans, finite words, tuples,
vectors, matrices, functions and streams have separate type constructors.
Dimensioned scalars also carry a seven-component rational SI exponent vector;
addition requires equal dimensions, multiplication adds dimensions, and
transcendental functions require dimensionless arguments. Angle conventions and
coordinate frames are explicit semantic annotations, not inferred from an
untyped number. A framed vector's type carries its fully qualified nominal frame
identity; a framed matrix carries source and target identities. Each identity
contains namespace, major/minor version and label. The full annotation is
serialized and participates in sharing. A frame change requires an explicit
typed template and its coordinate-realization obligations, not a relabelled
vector. Bare mathematical vectors carry no implied physical frame.

An integer is $(s,a_0,\ldots,a_{k-1})$ with $s\in\B$ and
\[
 z=(-1)^s\sum_{i=0}^{k-1}a_i\beta^i,
 \qquad \beta=2^{64},\qquad 0\le a_i<\beta.
\]
Zero has $s=0,k=0$; otherwise $a_{k-1}\ne0$. There is no format-level bound on
$k$. A resource-limited implementation must return an explicit resource result,
never silently discard high limbs. A rational is $(p,q)$ with $q>0$,
$\gcd(|p|,q)=1$; its zero is $(0,1)$.

An algebraic real has a primitive irreducible polynomial $P\in\Z[X]$ with
positive leading coefficient and an index selecting one of its distinct real
roots in increasing order. This fixes its value and canonical mathematical
identity. A rational isolating interval is a separate witness, not part of that
identity. Checking irreducibility, isolating roots and obtaining a canonical
minimal polynomial are explicit exact-algebra obligations. A general operator
term has canonical \emph{syntax}; different syntax may denote the same real.

\subsection{Exact arithmetic on finite machine words}
The literal full adder takes three bits $a,b,c\in\B$ and returns
\[
 s=a\oplus b\oplus c,\qquad
 c'= (a\land b)\lor\bigl(c\land(a\oplus b)\bigr),
 \qquad a+b+c=s+2c'.
\]
For subtraction with incoming borrow $d$, its bit and outgoing borrow are
\[
 r=a\oplus b\oplus d,\qquad
 d'=(\neg a\land b)\lor\bigl(d\land\neg(a\oplus b)\bigr),
 \qquad a-b-d=r-2d'.
\]
Feed the outgoing carry or borrow of position $i$ to position $i+1$, starting
with zero. Signed values are handled by separate sign and magnitude, not by
discarding an infinite sign extension. A product begins with partial bits
$p_{i,j}=a_i\land b_j$ in position $i+j$; finitely many full-adder chains sum
these shifted rows, retaining the final carry. Exact magnitude comparison scans
from the most significant padded bit, with equality flag $e=1$ and greater flag
$g=0$, updating simultaneously
\[
 g'=g\lor(e\land a\land\neg b),\qquad
 e'=e\land\neg(a\oplus b).
\]
This comparison supplies the branch in division; no approximate sign oracle is
needed for integers. The bit equations also establish the limb implementation:
$a+b=(a\oplus b)+2(a\land b)$ and the two carry terms are disjoint, giving
the full-adder identity by substitution. Summing the bit identities with powers
of two telescopes all internal carries and borrows.

Limb storage does not force modular arithmetic. For nonnegative addition, pad
the shorter operand with zero limbs, set $c_0=0$, and define
\begin{align*}
 u_i&=a_i+b_i+c_i,&
 r_i&=u_i-\beta\left\lfloor u_i/\beta\right\rfloor,&
 c_{i+1}&=\left\lfloor u_i/\beta\right\rfloor.
\end{align*}
Append the final nonzero carry. A native implementation may realize $u_i$ with
two-word operations or bitwise carry logic; it cannot assume that $u_i$ fits one
limb. Subtraction with $A\ge B$ uses $c_0=0$ and
\[
 u_i=a_i-b_i-c_i,\qquad
 c_{i+1}=[u_i<0],\qquad r_i=u_i+\beta c_{i+1}.
\]
Signed addition selects addition or magnitude subtraction according to signs
and exact magnitude order, then normalizes zero.

Multiplication is finite convolution followed by carrying, starting with $c_0=0$:
\[
 v_j=c_j+\sum_{i+\ell=j}a_i b_\ell,
 \qquad r_j=v_j\bmod\beta,
 \qquad c_{j+1}=\lfloor v_j/\beta\rfloor.
\]
Continue carrying until it vanishes. The accumulator $v_j$ is itself an exact
multiword quantity; using one overflowing hardware word would violate the rule.
For division by $D>0$, process dividend bits from most to least significant:
\[
 Q_0=R_0=0,\quad u_i=2R_i+b_i,\quad e_i=[u_i\ge D],\quad
 R_{i+1}=u_i-e_iD,\quad Q_{i+1}=2Q_i+e_i.
\]
The dividend prefix processed after step $i$ is $DQ_i+R_i$, with
$0\le R_i<D$. Signed Euclidean division is the unique
$z=dq+r$ with $d\ne0$ and $0\le r<|d|$; sign conversion must preserve this
convention, including negative dividends. Integer powers with nonnegative
exponent use finite repeated multiplication; negative powers produce a rational
only when the base is nonzero.

\begin{proposition}[No mathematical fixed-width overflow]
For finite integer operands, the preceding finite procedures return their exact
sum, difference, product, or Euclidean quotient and remainder, provided every
intermediate carry is retained and sufficient storage is available.
\end{proposition}
\begin{proof}
For addition and subtraction, multiply each defining equation by $\beta^i$
and sum: the carry terms telescope, leaving the represented operand relation.
Convolution represents the product by distributivity, and carrying preserves
each coefficient's value. For division, the prefix invariant follows by
substitution, while $0\le R_i<D$ implies $0\le2R_i+b_i<2D$ and hence preserves
the remainder bound. These arguments are conditional on retaining every term;
they establish no bound on time or allocated memory.
\end{proof}

\subsection{Positive roots and elementary functions}
\code{SQRT}$(x)$ denotes the unique $y\ge0$ satisfying $y^2=x$, with domain
$x\ge0$. \code{ROOT}$(n,x)$ requires integer $n\ge1$: for even $n$, it is the
nonnegative root and requires $x\ge0$; for odd $n$, it is the unique real root.
An arbitrary polynomial root uses the algebraic root index above, or an explicit
existence-and-uniqueness predicate for a real-term root. No unnamed branch is
selected. Thus $\sqrt{x^2}=|x|$ for real $x$, and cancellation to $x$ requires
$x\ge0$. Complex branches are outside this real profile.

The real elementary declarations have ordinary mathematical denotations:
\begin{align*}
 \exp x&=\sum_{n=0}^{\infty}\frac{x^n}{n!},&
 \sin x&=\sum_{n=0}^{\infty}\frac{(-1)^n x^{2n+1}}{(2n+1)!},\\
 \cos x&=\sum_{n=0}^{\infty}\frac{(-1)^n x^{2n}}{(2n)!},&
 \log x&=\int_1^x\frac{du}{u}\quad(x>0),\\
 \operatorname{atan}x&=\int_0^x\frac{du}{1+u^2},&
 \pi&=4\operatorname{atan}(1).
\end{align*}
These definitions permit finite operator terms for infinite-process values.
They are not a request to store an infinite series or an assertion that a finite
limb evaluator can return all digits. A certified finite approximation needs
separate truncation and rounding bounds. A raw approximate evaluation cannot
serve as an exact equality proof.

\code{ATAN2}$(y,x)$ is undefined at $(0,0)$ and uses range $(-\pi,\pi]$:
\[
 \operatorname{atan2}(y,x)=
 \begin{cases}
 \operatorname{atan}(y/x),&x>0,\\
 \operatorname{atan}(y/x)+\pi,&x<0,\ y\ge0,\\
 \operatorname{atan}(y/x)-\pi,&x<0,\ y<0,\\
 \pi/2,&x=0,\ y>0,\\
 -\pi/2,&x=0,\ y<0.
 \end{cases}
\]
The literal source \code{OTAN2} remains a distinct named template with its own
equation and domain; it is never an alias silently replaced by this operator.
$\tan x$ requires $\cos x\ne0$; inverse sine and cosine use ranges
$[-\pi/2,\pi/2]$ and $[0,\pi]$ and domain $[-1,1]$. Floor is the unique
integer $n$ with $n\le x<n+1$; wrapping is
$x-P\lfloor x/P\rfloor$ for $P>0$, with the integer winding retained separately.

\subsection{Decisions, partiality and honest result states}
Exact rational order is decided by integer cross multiplication with positive
denominators. Algebraic order and equality can be decided by exact polynomial
arithmetic, common-root analysis and real-root isolation. For unequal algebraic
roots, sufficiently refined disjoint rational isolating intervals determine the
order; common roots are handled algebraically instead of waiting forever for
interval separation. Supplying those algorithms is an implementation obligation.

There is no general equality or sign oracle for unrestricted transcendental
terms. \code{COMPARE} returns one of \code{LT}, \code{EQ}, \code{GT}, or
\code{UNKNOWN}, each decided result carrying a valid justification.
\code{FLOOR}, branch selection and nonzero checks may likewise remain unresolved.
Outcomes distinguish a value or a retained term, an unmet domain obligation,
an unknown decision, and exhausted resources. These are not numeric zero.
\code{IF} is lazy: only the selected branch must be defined once its predicate
is proved; an undecided predicate preserves the conditional term.

An exact finite decimal input denotes a rational with its declared decimal
scale. A measurement with error is a value plus an explicit uncertainty model;
giving its central value an exact rational spelling removes no measurement
error. This distinction applies equally to station coordinates, timing and force
coefficients in an orbital seed.

\subsection{Scoped operator graphs and named dependencies}
An operator seed contains finite typed records. A node is
\[
 (\text{opcode},\ \text{result type},\ \text{ordered operand references},
 \ \text{opcode payload}).
\]
The fixed \code{XOP-R1} catalog in \path{formal/OPERATOR_CATALOG.md} supplies
each opcode's arity, type relation, denotation, domain and payload fields.
An unknown opcode is rejected; a user-defined formula is a named template built
from catalog operators, with an explicit interface and body. There is no implicit
``evaluate arbitrary text'' opcode or undeclared exact solver.

Each scope has an ordered node table and ordered exported references. Local
node operands refer only to earlier nodes. Lambda, finite sum/product, integral,
derivative and solution declarations carry nested scopes. Bound variables use
de Bruijn indices: zero denotes the nearest binder, and larger indices count
outward through lexical binders. The type environment fixes every index.
Substitution shifts indices when passing a binder and is capture-avoiding.
Nested scopes cannot reference an outer local node by an accidental numeric ID;
captured values enter through explicit binders. Names are display metadata;
renaming them does not change variable binding.

References have distinct tags for a local node, bound variable, old state slot,
current sample value, dependency export, raw sample timestamp, or converted GPST
timestamp. A template application additionally
identifies its template and supplies all parameters in declared order. Dependency
records carry the complete canonical seed/template bytes in the self-contained
form. A digest is a lookup aid; acceptance requires the expected complete bytes
and interface. A dependency-only manifest is explicitly an open package until
all required bytes are supplied. Dependency and template-call graphs are finite
and acyclic in this profile. Finite summation and delayed state do not imply
unrestricted recursive template calls.

Sample declarations identify their type, absolute time convention, provenance
and binding mode. A time convention has a versioned semantic profile identity;
the registered GPST convention supplies its stated scale and epoch, while a
custom convention includes a complete dependency/template mapping and its domain.
A name alone does not define a clock, resolve a leap second or prove monotonicity.
An inline sample is an indexed literal accompanied by its
absolute timestamp; an external sample is a declared input requirement.
\code{SAMPLE} selects the value bound to the current logical step, with its
timestamp retained. Missing data is unresolved input, not a guessed next sample.
Raw-time reference tag 5 returns the sample's exact rational in its declared
source-clock coordinate. Converted-time tag 6 explicitly applies the versioned
mapping and returns GPST SI seconds; the value and both timestamp references
read one immutable sample snapshot. Integration arguments are explicit
expressions such as $t=\tau_{\rm GPST}-\tau_{\rm epoch}$ and
$h=\tau_{\rm GPST}-\tau_{\rm previous}$, with matching time units and epoch
convention. A custom raw clock tick is not implicitly a second, and a missing
mapping cannot be replaced by an assumed offset.
A modular phase, spatial key, hash or template identifier does not supply the
state, timestamp or sample bytes it refers to.

\subsection{Canonical variable-length wire contract}
The canonical semantic stream begins with the eight ASCII octets
\code{XOPSEED1}, followed by profile numbers $(1,0)$ and the ordered tables
specified in the catalog. Every natural number uses minimal unsigned base-128
encoding: digit payloads occupy seven low bits, the high bit marks continuation,
the least significant digit comes first, and the final high bit is zero.
Zero is one zero octet; a multi-octet number cannot have a zero final payload.
Thus counts and references are not limited to a fixed machine word.

A byte string has a natural-number length followed by exactly that many octets;
a sequence has an element count followed by exactly that many schema-defined
elements. An integer has a sign octet, a limb count and that many little-endian
64-bit limbs. Rational and algebraic literals use the canonical value forms
already defined. Strings used as identifiers are restricted ASCII byte strings;
free-form provenance is opaque bytes and does not participate in name lookup.
Types, references and opcode payloads are encoded by their declared tagged
schemas. Nonminimal counts, redundant leading limbs, trailing bytes, nonzero
padding, invalid references and payload fields outside the selected schema have
no canonical decoding.

Canonical here means one encoding for the same ordered, scoped syntactic
record. It does not mean that all mathematically equal expressions, all sharing
patterns, or all valid topological node orders receive one universal byte string.
Structural sharing is permitted only for identical typed syntax under the same
bindings and immutable dependency identities. The strategy history records any
rewrite that changes the syntax while preserving its meaning.

\begin{proposition}[Unique literal and record decoding]
Given the catalog, a finite well-formed canonical stream has one decoded record
and its canonical re-encoding returns the original stream.
\end{proposition}
\begin{proof}
The continuation terminator uniquely delimits each minimal natural-number
encoding. Lengths and fixed-width limbs uniquely delimit strings and literals.
The record schema fixes field order and the selected tag fixes each variant.
Induction over finite sequences and nested scopes therefore determines every
field and its boundary. Rejecting alternative zero forms, redundant lengths and
trailing bytes makes re-encoding identical. This is a structural proof of the
specified grammar, not evidence about a parser implementation.
\end{proof}

To pack literal one-bit planes, let $C$ be the canonical octet stream, of length
$L$. Pad its last block with zero octets to a multiple of 512 octets. Interpret
each block as 64 little-endian words $w_0,\ldots,w_{63}$ and set
\[
 a_{r,c}=\bigl((w_r\gg c)\mathbin{\&}1\bigr),\qquad
 p_c=\sum_{r=0}^{63}a_{r,c}2^r.
\]
The envelope is eight ASCII octets \code{XOPPLAN1}, then the minimal encoding of
$L$, then the plane words in increasing block and plane order. Its size fixes
the required block count; after inverse transpose, padding must be zero and the
first $L$ octets must decode as \code{XOPSEED1}. Empty input has no seed grammar,
although the underlying byte permutation can act on an empty byte string.

\begin{proposition}[Bit-plane involution and information accounting]
The padded block transform $T$ satisfies $T(T(W))=W$. Including $L$ makes packing
injective on canonical seeds. The transform changes organization, not information
content or the number of bits required by its uncompressed padded blocks.
\end{proposition}
\begin{proof}
The output bit at $(c,r)$ is the input bit at $(r,c)$. Transposing again restores
every coordinate. The explicit length removes padding ambiguity. Consequently
distinct canonical inputs cannot share an envelope. A bijective permutation
cannot itself identify multiple arbitrary inputs with a smaller code space.
\end{proof}
This envelope is newly named. It neither changes nor retroactively reinterprets
R10's uint64/binary64 codec. A conversion from R10 must declare whether it
preserves raw bits, converts finite values to dyadic rationals, or replaces an
approximation by a separately justified symbolic model.

\subsection{Causal recurrence and the literal two-mask kernel}
The state table declares each slot's type, initial expression and next-state
expression. Reading slot $j$ at logical step $t$ is precisely
\[
 s_j(0)=\operatorname{init}_j,\qquad
 s_j(t+1)=E_j(s(t),d(t));
 \qquad s_j=\operatorname{DELAY}(\operatorname{init}_j,E_j).
\]
All right-hand sides read the same old state and the same input snapshot. The
new tuple commits simultaneously after all required domains and decisions are
resolved. A self-loop in an instantaneous operator DAG is inadmissible: every
recurrence cycle must cross a declared \code{DELAY}. An algebraic fixed-point
equation is a separate solution declaration needing existence, uniqueness
and branch obligations; it is not an execution order.

For $w$-bit masks, define the retained whole-word operation
\begin{align*}
 a&=x\land A\land V,& n&=a\land N,& h&=\pop(n\land B),\\
 \operatorname{ASA}(x;V,A,N,B)&=
 \begin{cases}0,&h>0,\\n,&h=0.\end{cases}
\end{align*}
The two ordered sets and literal J--K update are
\begin{align*}
 x_t&=E_X(q_t,d_t),&
 y_t^{A}&=\operatorname{ASA}(x_t;V_A,A_A,N_A,B_A),\\
 y_t&=\operatorname{ASA}(y_t^{A};V_B,A_B,N_B,B_B),\\
 J_t&=E_J(q_t,y_t,d_t),& K_t&=E_K(q_t,y_t,d_t),\\
 q_{t+1}&=\bigl[(J_t\land\neg_w q_t)
       \lor(\neg_w K_t\land q_t)\bigr]\land V_q.
\end{align*}
Here $\neg_w q=(2^w-1)\mathbin{\oplus}q$ and every expression has
$\mathsf{Bits}(w)$ type. Both intermediate $(a,n,h,y)$ tuples are observables.
Absorption zeroes the whole output of its stage; it neither clears selected bits
individually nor silently substitutes a hold transition. The mathematical width
$w$ is explicit; the retained executable orbital kernel still uses its specified
32-bit word profile.

\begin{proposition}[Causal uniqueness and conditional finite periodicity]
For fixed initial states and an input stream, a typed finite instantaneous DAG
whose required operations are defined determines at most one next state at each
step. If all operations are total on the reachable states, it determines a unique
state stream. A fixed autonomous transition on $m$ active Boolean state bits is
eventually periodic, with preperiod plus period at most $2^m$.
\end{proposition}
\begin{proof}
Topological induction determines every demanded instantaneous node on the
admitted branches from old-state and input values; simultaneous commitment
then determines the next tuple. Induction
on steps gives the stream. For the finite autonomous case, among $2^m+1$ states
two coincide; determinism makes their subsequent tails coincide. Changing the
input, strategy, timestamp, unbounded integer state or physical continuous state
does not meet that finite-autonomous hypothesis.
\end{proof}

\subsection{Vectors, calculus and orbital expressions}
Vector and matrix operators are typed compositions of scalar arithmetic:
$u\cdot v=\sum_i u_i v_i$, $\|u\|=\sqrt{u\cdot u}$, and
$(AB)_{ij}=\sum_k A_{ik}B_{kj}$. Cross product is restricted to three
components. A normalized vector requires nonzero norm; a matrix inverse requires
a square matrix and nonzero determinant. Polynomial, Chebyshev, Legendre,
finite Fourier and rational expressions expand to declared finite templates.
The catalog includes finite sum/product, differentiation, gradient and Jacobian
binding, definite integration and an initial-value solution declaration.

\code{DIFF}$(f,x)$ denotes a derivative only where the stated limit exists.
\code{INTEGRAL}$(f,a,b)$ denotes an oriented Riemann integral when the finite
interval integrand is integrable; an improper integral has a separate limit
declaration. \code{ODE}$(F,t_0,s_0,t,I)$ denotes the value at $t\in I$ of the
unique continuous state satisfying
\[
 s(u)=s_0+\int_{t_0}^{u}F(v,s(v))\,dv\qquad(u\in I),
\]
with a Riemann-integrable right-hand side along that state. On continuous
pieces the solution is differentiable; a coefficient jump need not admit a
derivative at the junction. The named continuous \code{FLOW} uses this
integral-equation meaning. State may be a tuple of position, velocity and other
components with different dimensions, including explicitly framed vectors.
Each derivative component retains its state's nominal frame and has its state
dimension divided by time. Rotating-frame terms must occur in $F$; the type
annotation does not invent them. A local existence theorem or local regularity
condition alone is not a certificate of a global interval.
If existence, uniqueness or continuation is unresolved, the node remains a
declaration with obligations, rather than an evaluated exact orbit.

A finite RK4 template can be represented exactly as expressions, for example
\begin{align*}
 k_1&=F(t,s),&k_2&=F(t+h/2,s+hk_1/2),\\
 k_3&=F(t+h/2,s+hk_2/2),&k_4&=F(t+h,s+hk_3),\\
 \operatorname{RK4}_h(t,s)&=s+\frac{h}{6}(k_1+2k_2+2k_3+k_4).
\end{align*}
Exact evaluation of this finite expression does not make it equal to the
continuous ODE solution. Discretization, uncertain initial conditions, unmodelled
forces and discontinuous external events remain separate mathematical or
physical questions. All force and station expressions must retain their units,
frames, absolute time and explicit singular domains. A finite symbolic term is
not an error bound on its physical denotation.

\subsection{Proof rules and strategy self-improvement}
Let $\Gamma$ contain typed assumptions and proven domain predicates. A permitted
rewrite carries the judgement
\[
 \Gamma\vdash e\simeq e':\tau,
 \qquad
 \operatorname{Dom}_{\Gamma}(e)=\operatorname{Dom}_{\Gamma}(e'),
 \quad \llbracket e\rrbracket=\llbracket e'\rrbracket
 \text{ on that domain}.
\]
Domain equality is essential: simplifying $x/x$ to $1$ without $x\ne0$ changes
the defined inputs. The proof vocabulary contains assumption, reflexivity,
symmetry, transitivity, typed congruence, capture-avoiding substitution,
definition expansion and explicit ordered-field/domain lemmas. Branch rules
require their predicates; differentiation and integration rewrites additionally
require their analytic hypotheses. A rule name or a source digest alone is not
a proof of its side conditions.

Structural common-subexpression sharing, exact constant folding and eliminating
a pure unused node preserve the denotation of the demanded exports when the
typing, scope and domain obligations are satisfied. Reassociation is justified
in exact integer/rational arithmetic but is not silently applied to R10 binary64
execution. Eliminating an observed ASA trace, changing a sample read, or changing
query cadence is not observationally neutral merely because a final word matches.

For stateful replacement use a relation $\mathcal R$ between old and new states.
Its proof must relate the initial states and, for every admitted input, preserve
definedness, all declared observations and the relation after one synchronous
step. Induction then gives equal observable streams. Absolute timestamps and
selected sample identities belong to the observations whenever a query schedule
is exported; a change in those requires a separately named strategy contract.

A strategy version is an immutable tuple
\[
 S_v=(\text{rule catalog},\ \text{priority policy},\ \text{cost model},
       \ \text{proof system ID},\ \text{resource policy}).
\]
It may propose a new expression graph or a new strategy $S_{v+1}$. Acceptance
requires a finite proof object under the already accepted proof system, a
preserved input/output contract, and the stated termination/resource argument.
The proposal cannot justify itself by replacing its own checker or dropping an
unmet premise. A proof-system change creates a new explicit trust boundary.
Strategies and proposal histories may be unbounded, so finite Boolean
periodicity does not apply to the combined evolving system.

\begin{proposition}[Conditional preservation under strategy replacement]
Suppose the accepted proof rules are sound for the declared denotation, each
accepted replacement carries a valid domain-preserving proof or stream
simulation relation, and proof checking admits only valid derivations. Then any
finite chain of accepted replacements preserves its declared observations.
\end{proposition}
\begin{proof}
A sound proof preserves one replacement's domain and observations. Compose
these equalities, or the stepwise stream relations, along the finite chain.
Induction gives preservation after every accepted replacement. This establishes
a conditional theorem; it supplies neither a proof-checker implementation nor
an unconditional claim that proposed derivations are valid.
\end{proof}

Logical strategy improvement can now be stated precisely. Sharing an identical
pure subgraph reduces duplicated graph storage; memoizing it avoids a repeated
evaluation only when the complete binding, state step and dependency identity
are the same. A strictly decreasing natural-valued rewrite measure proves
termination of that rewrite phase; a finite proposal budget also terminates a
search without proving it finds a better expression. A predicted reduction in
operator count is a claim under the declared cost model, not a measured speedup
or a guarantee of lower big-integer time, memory use, physical error or global
optimality. These are the available reasoning-based improvements in this
subversion; execution and machine proof remain future obligations.
